\documentclass[11pt,a4paper]{article}

\usepackage[spanish]{babel}
\usepackage[utf8]{inputenc}
\usepackage[T1]{fontenc}
\usepackage{lmodern}
\usepackage{amsmath}
\usepackage{setspace}
\usepackage{geometry}
\usepackage{hyperref}

\geometry{margin=2.5cm}

\title{El mito del incentivo y la colonización capitalista de la motivación}
\author{Vicente Domínguez Arca}
\date{}

\begin{document}

\maketitle

Uno de los argumentos más recurrentes empleados por los defensores del neoliberalismo contemporáneo para desacreditar cualquier propuesta orientada hacia lo común consiste en apelar a la supuesta desaparición del incentivo. Según esta tesis, expresada frecuentemente con apariencia de evidencia autojustificada, los sistemas basados en la cooperación, la gestión común o la redistribución estarían condenados al fracaso porque los individuos perderían la motivación necesaria para actuar, producir, innovar o contribuir al conjunto social.

Sin embargo, esta afirmación presenta un problema epistemológico elemental: rara vez se acompaña de una fundamentación rigurosa sobre qué se entiende exactamente por incentivo. La tesis suele funcionar como un axioma encubierto. El incentivo aparece como una noción supuestamente evidente, de modo que cualquier petición de justificación es interpretada como una negación de una verdad autoevidente. Pero ninguna repetición sistemática transforma una afirmación en una demostración. La frecuencia con la que se enuncia una idea no aumenta su capacidad explicativa.

Resulta útil comenzar por una consideración etimológica. Incentivo deriva de diversas expresiones latinas relacionadas con la acción de encender o avivar el fuego. Metafóricamente, designa aquello que despierta, moviliza, excita o impulsa una conducta. El incentivo remite, por tanto, al estímulo.

Y aquí aparece una primera dificultad para la tesis neoliberal. Mientras permanezcamos despiertos y participemos en el mundo, nos encontramos sometidos a una estimulación constante. Los seres humanos son organismos abiertos que interactúan permanentemente con su entorno. Incluso en situaciones de privación sensorial parcial, otros estímulos continúan operando. La experiencia humana está constituida por una continua interacción con estímulos físicos, emocionales, cognitivos y sociales.

Desde una perspectiva ontológica elemental, el estímulo no es una excepción dentro de la vida humana: es una de sus condiciones constitutivas. Del mismo modo que no existe sociedad sin interacción, tampoco existe sujeto social completamente privado de estimulación. La propia existencia en comunidad implica una red continua de afecciones mutuas, expectativas, necesidades, proyectos, afectos y responsabilidades.

Por ello, afirmar que un sistema basado en lo común carece de incentivos constituye una formulación conceptualmente problemática. Mientras existan personas, existirán estímulos. Mientras existan relaciones sociales, existirán motivaciones. El problema no es la existencia o inexistencia del incentivo, sino la forma concreta que adopta.

Lo que realmente suele querer decir el discurso neoliberal no es que desaparezcan los estímulos, sino que desaparecerían ciertos comportamientos específicos considerados deseables. El incentivo queda entonces redefinido de manera implícita. Ya no significa aquello que moviliza a las personas, sino aquello que las orienta hacia determinadas prácticas económicas compatibles con la reproducción ampliada del capital.

La operación conceptual es importante. Bajo esta redefinición, una persona sólo aparece como incentivada cuando su actividad desemboca en ciertos resultados previamente seleccionados: consumir determinados bienes, competir por determinadas posiciones, acumular patrimonio, maximizar rentabilidades o incrementar capacidades de adquisición dentro del mercado.

La cuestión deja entonces de ser ontológica para convertirse en ideológica.

No se discute si existen estímulos, sino cuáles son considerados legítimos. No se analiza qué motiva a las personas, sino qué motivaciones resultan funcionales para una estructura económica concreta.

Este desplazamiento resulta especialmente problemático porque gran parte de los deseos individuales no surgen espontáneamente en el vacío. Las sociedades modernas dedican enormes recursos materiales a producir sistemas de estimulación orientados. La publicidad, el marketing, la construcción de prestigio simbólico, la industria cultural y los mecanismos de diferenciación social participan activamente en la configuración de aquello que aparece como deseable.

Lo que los individuos desean no es independiente de los sistemas que producen significados. El deseo es también un fenómeno social.

Por ello, cuando se afirma que un sistema basado en lo común destruiría los incentivos, muchas veces lo que realmente se está afirmando es que dejaría de garantizar determinados patrones de deseo orientados al consumo y a la acumulación privada. El incentivo se identifica entonces con la reproducción de la lógica mercantil.

Pero esta identificación carece de necesidad lógica.

Las personas pueden sentirse movilizadas por el reconocimiento, el afecto, la cooperación, la creatividad, el prestigio intelectual, la participación política, la pertenencia comunitaria, la curiosidad científica, la vocación artística, la solidaridad o el simple deseo de contribuir al bienestar colectivo. La historia humana ofrece innumerables ejemplos de actividades sostenidas por motivaciones distintas de la acumulación privada.

Esto obliga a introducir una distinción conceptual más precisa entre incentivo y premio.

El incentivo designa la existencia de estímulos capaces de movilizar conductas. En este sentido, es una constante antropológica difícilmente eliminable. El premio, en cambio, remite al resultado esperado de una determinada acción.

Los sistemas sociales distribuyen premios de formas diferentes.

El capitalismo tiende a organizar los premios en torno a la apropiación diferencial de recursos, riqueza, propiedad y capacidad de consumo. Pero esta no constituye la única forma posible de estructurar recompensas.

En una organización social basada en lo común, el premio puede adoptar una forma distinta. Puede consistir en la expectativa razonable de que la mejora del sistema repercutirá sobre la propia vida. Allí donde existen garantías colectivas robustas, el bienestar individual deja de depender exclusivamente de la capacidad competitiva de cada sujeto y pasa a estar vinculado a la evolución favorable del conjunto.

La recompensa ya no aparece como apropiación diferencial, sino como participación creciente en un bienestar compartido.

Dicho de otro modo, el premio emerge de la propia estabilidad del sistema.

Cuanto más capaz sea una sociedad de garantizar vivienda, educación, salud, tiempo libre, participación política, seguridad material y desarrollo personal para todos sus miembros, mayor será el interés racional de cada individuo en contribuir a la reproducción de ese sistema.

Desde esta perspectiva, el problema no es la ausencia de incentivos, sino la credibilidad de las garantías.

Y aquí aparece una dificultad fundamental para los sistemas basados en la acumulación privada. Allí donde existen gradientes estructurales de acceso, capacidad decisional y libertad efectiva, las garantías dejan de ser universales y se convierten en condiciones particulares. Lo que se presenta como derecho termina subordinado a la posición ocupada dentro de la estructura económica.

La promesa de bienestar se transforma entonces en una promesa condicional.

Esta es precisamente una de las formas más profundas de restricción de la libertad. No porque desaparezcan las opciones formales, sino porque las posibilidades efectivas quedan distribuidas de manera desigual.

La paradoja final es que quienes suelen presentar el capitalismo como el sistema de los incentivos confunden sistemáticamente motivación humana con estímulo mercantil. Pero los seres humanos no necesitan el mercado para desear, crear, cooperar o proyectar futuros. Lo que necesitan es un horizonte de sentido y una expectativa razonable de que su participación en el tejido social contribuirá también a mejorar sus propias condiciones de existencia.

La defensa neoliberal del incentivo no describe una necesidad antropológica universal. Describe, más modestamente, la necesidad funcional de un sistema económico que requiere convertir determinadas formas de deseo en mecanismos permanentes de reproducción de la acumulación.

\end{document}

